\documentclass[11pt]{article}
\usepackage[a4paper,margin=27mm]{geometry}
\usepackage{amsmath,amssymb,mathtools}
\usepackage{fontspec,unicode-math}
\setmainfont{Latin Modern Roman}
\setmathfont{Latin Modern Math}
\usepackage{xurl,hyperref}
\hypersetup{colorlinks=true,urlcolor=blue,linkcolor=blue}
\newcommand{\sref}[2]{\hyperlink{source:#1:#2}{[S#2]}}
\newcommand{\cataloguescope}[1]{\paragraph{Mathematical question.} #1}
\begin{document}
\section{Conjecture 3.24 — Bertrand}\label{statement}

\subsection{Definitions and mathematical statement}

\paragraph{Shared notation.}$\mathbb N=\{1,2,\ldots\}$, and $\mathbb Z,\mathbb Q,\mathbb R,\mathbb C$ denote the integers, rationals, reals and complex numbers. Any other conventions are specified locally.


For $0<|q|<1$, put $E_4(q)=1+240\sum_{n\ge1}n^3q^n/(1-q^n)$ and $\Delta(q)=q\prod_{n\ge1}(1-q^n)^{24}$, and define the classical modular invariant $J(q)=E_4(q)^3/\Delta(q)$. Set $D=q\,d/dq$. Two nonzero complex numbers are multiplicatively dependent if $q_1^a q_2^b=1$ for integers $(a,b)\ne(0,0)$.
\paragraph{Statement recorded in the supplied source.}
For every $n\ge1$ and algebraic $q_1,\ldots,q_n$ with $0<|q_i|<1$, does algebraic dependence over $\mathbb Q$ of the $3n$ numbers $J(q_i),DJ(q_i),D^2J(q_i)$ imply that some distinct $i,j$ have multiplicatively dependent $q_i,q_j$? \sref{statement}{2}
\subsection{Short English statement}
Must any algebraic relation among modular invariant values and their first two logarithmic derivatives at algebraic arguments force a multiplicative relation between two arguments?
\subsection{Sources}
\begin{description}
\item[\hypertarget{source:statement:1}{S1}] ULAM.AI and the UnsolvedMath Contributors, UnsolvedMath: A Curated Collection of Open Mathematics Problems, catalogue entry. CC BY 4.0. \url{https://huggingface.co/datasets/ulamai/UnsolvedMath}.
\item[\hypertarget{source:statement:2}{S2}] Michel Waldschmidt, Open Diophantine Problems, arXiv:math/0312440v2 (24 January 2004), Conjecture 3.24. \url{https://arxiv.org/abs/math/0312440}.
\end{description}
\label{end:statement}
\end{document}
